\documentclass[12pt, a4paper]{article}

\usepackage[a4paper, margin=1in]{geometry}

\usepackage[utf8]{inputenc}
\usepackage[mathscr]{euscript}
\usepackage{amssymb,amsmath,amsthm,amsfonts}
\usepackage[shortlabels]{enumitem}
\usepackage{multicol,multirow}
\usepackage{lipsum}
\usepackage{balance}
\usepackage{calc}
\usepackage[colorlinks=true,citecolor=blue,linkcolor=blue]{hyperref}
\usepackage{import}
\usepackage{xifthen}
\usepackage{pdfpages}
\usepackage{transparent}
\usepackage{listings}

\newcommand{\incfig}[2][1.0]{
    \def\svgwidth{#1\columnwidth}
    \import{./figures/}{#2.pdf_tex}
}

\newcommand{\incimg}[2][1.0]{
  \includegraphics[width=#1\columnwidth]{./figures/#2}
}

\newlist{enumproof}{enumerate}{4}
\setlist[enumproof,1]{label=\arabic*., parsep=1em}
\setlist[enumproof,2]{label=\arabic{enumproofi}.\arabic*., parsep=1em}
\setlist[enumproof,3]{label=\arabic{enumproofi}.\arabic{enumproofii}.\arabic*., parsep=1em}
\setlist[enumproof,4]{label=\arabic{enumproofi}.\arabic{enumproofii}.\arabic{enumproofiii}.\arabic*., parsep=1em}

\renewcommand{\qedsymbol}{\ensuremath{\blacksquare}}

\lstdefinestyle{mystyle}{
  language=Python, % Set the language to Python
  commentstyle=\color{codegray}\itshape, % Python comments (#)
  keywordstyle=\color{orange}\bfseries,  % Python reserved keywords
  stringstyle=\color{mauve},             % Strings
  basicstyle={\ttfamily\footnotesize},   % Monospace font
  breakatwhitespace=true,         
  breaklines=true,                 
  captionpos=b,                    
  keepspaces=true,                 
  numbers=none,                    
  tabsize=2,
  showstringspaces=false,
  % Extra keyword groups for Python
  morekeywords=[2]{self, cls, __init__, __str__, __repr__}, 
  keywordstyle=[2]\color{identblue},     % Special identifiers / methods
  morekeywords=[3]{int, float, bool, str, list, dict, set, tuple, object, None, True, False},
  keywordstyle=[3]\color{codegreen},     % Built-in types
  morekeywords=[4]{print, input, open, range, len, type, isinstance, super, enumerate, zip, map, filter, sum, min, max, sorted},
  keywordstyle=[4]\color{cyan},          % Built-in functions
}

% Things Lie
\newcommand{\kb}{\mathfrak b}
\newcommand{\kg}{\mathfrak g}
\newcommand{\kh}{\mathfrak h}
\newcommand{\kn}{\mathfrak n}
\newcommand{\ku}{\mathfrak u}
\newcommand{\kz}{\mathfrak z}
\DeclareMathOperator{\Ext}{Ext} % Ext functor
\DeclareMathOperator{\Tor}{Tor} % Tor functor
\newcommand{\gl}{\opname{\mathfrak{gl}}} % frak gl group
\renewcommand{\sl}{\opname{\mathfrak{sl}}} % frak sl group chktex 6

% More script letters etc.
\newcommand{\SA}{\mathcal A}
\newcommand{\SB}{\mathcal B}
\newcommand{\SC}{\mathcal C}
\newcommand{\SF}{\mathcal F}
\newcommand{\SG}{\mathcal G}
\newcommand{\SH}{\mathcal H}
\newcommand{\OO}{\mathcal O}

\newcommand{\SCA}{\mathscr A}
\newcommand{\SCB}{\mathscr B}
\newcommand{\SCC}{\mathscr C}
\newcommand{\SCD}{\mathscr D}
\newcommand{\SCE}{\mathscr E}
\newcommand{\SCF}{\mathscr F}
\newcommand{\SCG}{\mathscr G}
\newcommand{\SCH}{\mathscr H}

% Mathfrak primes
\newcommand{\km}{\mathfrak m}
\newcommand{\kp}{\mathfrak p}
\newcommand{\kq}{\mathfrak q}

% number sets
\newcommand{\RR}[1][]{\ensuremath{\ifstrempty{#1}{\mathbb{R}}{\mathbb{R}^{#1}}}}
\newcommand{\NN}[1][]{\ensuremath{\ifstrempty{#1}{\mathbb{N}}{\mathbb{N}^{#1}}}}
\newcommand{\ZZ}[1][]{\ensuremath{\ifstrempty{#1}{\mathbb{Z}}{\mathbb{Z}^{#1}}}}
\newcommand{\QQ}[1][]{\ensuremath{\ifstrempty{#1}{\mathbb{Q}}{\mathbb{Q}^{#1}}}}
\newcommand{\CC}[1][]{\ensuremath{\ifstrempty{#1}{\mathbb{C}}{\mathbb{C}^{#1}}}}
\newcommand{\PP}[1][]{\ensuremath{\ifstrempty{#1}{\mathbb{P}}{\mathbb{P}^{#1}}}}
\newcommand{\HH}[1][]{\ensuremath{\ifstrempty{#1}{\mathbb{H}}{\mathbb{H}^{#1}}}}
\newcommand{\FF}[1][]{\ensuremath{\ifstrempty{#1}{\mathbb{F}}{\mathbb{F}^{#1}}}}
% expected value
\newcommand{\EE}{\ensuremath{\mathbb{E}}}
\newcommand{\charin}{\text{ char }}
\DeclareMathOperator{\sign}{sign}
\DeclareMathOperator{\Aut}{Aut}
\DeclareMathOperator{\Inn}{Inn}
\DeclareMathOperator{\Syl}{Syl}
\DeclareMathOperator{\Gal}{Gal}
\DeclareMathOperator{\GL}{GL} % General linear group
\DeclareMathOperator{\SL}{SL} % Special linear group

%---------------------------------------
% BlackBoard Math Fonts :-
%---------------------------------------

%Captital Letters
\newcommand{\bbA}{\mathbb{A}}	\newcommand{\bbB}{\mathbb{B}}
\newcommand{\bbC}{\mathbb{C}}	\newcommand{\bbD}{\mathbb{D}}
\newcommand{\bbE}{\mathbb{E}}	\newcommand{\bbF}{\mathbb{F}}
\newcommand{\bbG}{\mathbb{G}}	\newcommand{\bbH}{\mathbb{H}}
\newcommand{\bbI}{\mathbb{I}}	\newcommand{\bbJ}{\mathbb{J}}
\newcommand{\bbK}{\mathbb{K}}	\newcommand{\bbL}{\mathbb{L}}
\newcommand{\bbM}{\mathbb{M}}	\newcommand{\bbN}{\mathbb{N}}
\newcommand{\bbO}{\mathbb{O}}	\newcommand{\bbP}{\mathbb{P}}
\newcommand{\bbQ}{\mathbb{Q}}	\newcommand{\bbR}{\mathbb{R}}
\newcommand{\bbS}{\mathbb{S}}	\newcommand{\bbT}{\mathbb{T}}
\newcommand{\bbU}{\mathbb{U}}	\newcommand{\bbV}{\mathbb{V}}
\newcommand{\bbW}{\mathbb{W}}	\newcommand{\bbX}{\mathbb{X}}
\newcommand{\bbY}{\mathbb{Y}}	\newcommand{\bbZ}{\mathbb{Z}}

%---------------------------------------
% MathCal Fonts :-
%---------------------------------------

%Captital Letters
\newcommand{\mcA}{\mathcal{A}}	\newcommand{\mcB}{\mathcal{B}}
\newcommand{\mcC}{\mathcal{C}}	\newcommand{\mcD}{\mathcal{D}}
\newcommand{\mcE}{\mathcal{E}}	\newcommand{\mcF}{\mathcal{F}}
\newcommand{\mcG}{\mathcal{G}}	\newcommand{\mcH}{\mathcal{H}}
\newcommand{\mcI}{\mathcal{I}}	\newcommand{\mcJ}{\mathcal{J}}
\newcommand{\mcK}{\mathcal{K}}	\newcommand{\mcL}{\mathcal{L}}
\newcommand{\mcM}{\mathcal{M}}	\newcommand{\mcN}{\mathcal{N}}
\newcommand{\mcO}{\mathcal{O}}	\newcommand{\mcP}{\mathcal{P}}
\newcommand{\mcQ}{\mathcal{Q}}	\newcommand{\mcR}{\mathcal{R}}
\newcommand{\mcS}{\mathcal{S}}	\newcommand{\mcT}{\mathcal{T}}
\newcommand{\mcU}{\mathcal{U}}	\newcommand{\mcV}{\mathcal{V}}
\newcommand{\mcW}{\mathcal{W}}	\newcommand{\mcX}{\mathcal{X}}
\newcommand{\mcY}{\mathcal{Y}}	\newcommand{\mcZ}{\mathcal{Z}}


%---------------------------------------
% Bold Math Fonts :-
%---------------------------------------

%Captital Letters
\newcommand{\bmA}{\boldsymbol{A}}	\newcommand{\bmB}{\boldsymbol{B}}
\newcommand{\bmC}{\boldsymbol{C}}	\newcommand{\bmD}{\boldsymbol{D}}
\newcommand{\bmE}{\boldsymbol{E}}	\newcommand{\bmF}{\boldsymbol{F}}
\newcommand{\bmG}{\boldsymbol{G}}	\newcommand{\bmH}{\boldsymbol{H}}
\newcommand{\bmI}{\boldsymbol{I}}	\newcommand{\bmJ}{\boldsymbol{J}}
\newcommand{\bmK}{\boldsymbol{K}}	\newcommand{\bmL}{\boldsymbol{L}}
\newcommand{\bmM}{\boldsymbol{M}}	\newcommand{\bmN}{\boldsymbol{N}}
\newcommand{\bmO}{\boldsymbol{O}}	\newcommand{\bmP}{\boldsymbol{P}}
\newcommand{\bmQ}{\boldsymbol{Q}}	\newcommand{\bmR}{\boldsymbol{R}}
\newcommand{\bmS}{\boldsymbol{S}}	\newcommand{\bmT}{\boldsymbol{T}}
\newcommand{\bmU}{\boldsymbol{U}}	\newcommand{\bmV}{\boldsymbol{V}}
\newcommand{\bmW}{\boldsymbol{W}}	\newcommand{\bmX}{\boldsymbol{X}}
\newcommand{\bmY}{\boldsymbol{Y}}	\newcommand{\bmZ}{\boldsymbol{Z}}
%Small Letters
\newcommand{\bma}{\boldsymbol{a}}	\newcommand{\bmb}{\boldsymbol{b}}
\newcommand{\bmc}{\boldsymbol{c}}	\newcommand{\bmd}{\boldsymbol{d}}
\newcommand{\bme}{\boldsymbol{e}}	\newcommand{\bmf}{\boldsymbol{f}}
\newcommand{\bmg}{\boldsymbol{g}}	\newcommand{\bmh}{\boldsymbol{h}}
\newcommand{\bmi}{\boldsymbol{i}}	\newcommand{\bmj}{\boldsymbol{j}}
\newcommand{\bmk}{\boldsymbol{k}}	\newcommand{\bml}{\boldsymbol{l}}
\newcommand{\bmm}{\boldsymbol{m}}	\newcommand{\bmn}{\boldsymbol{n}}
\newcommand{\bmo}{\boldsymbol{o}}	\newcommand{\bmp}{\boldsymbol{p}}
\newcommand{\bmq}{\boldsymbol{q}}	\newcommand{\bmr}{\boldsymbol{r}}
\newcommand{\bms}{\boldsymbol{s}}	\newcommand{\bmt}{\boldsymbol{t}}
\newcommand{\bmu}{\boldsymbol{u}}	\newcommand{\bmv}{\boldsymbol{v}}
\newcommand{\bmw}{\boldsymbol{w}}	\newcommand{\bmx}{\boldsymbol{x}}
\newcommand{\bmy}{\boldsymbol{y}}	\newcommand{\bmz}{\boldsymbol{z}}

%---------------------------------------
% Scr Math Fonts :-
%---------------------------------------

\newcommand{\sA}{{\mathscr{A}}}   \newcommand{\sB}{{\mathscr{B}}}
\newcommand{\sC}{{\mathscr{C}}}   \newcommand{\sD}{{\mathscr{D}}}
\newcommand{\sE}{{\mathscr{E}}}   \newcommand{\sF}{{\mathscr{F}}}
\newcommand{\sG}{{\mathscr{G}}}   \newcommand{\sH}{{\mathscr{H}}}
\newcommand{\sI}{{\mathscr{I}}}   \newcommand{\sJ}{{\mathscr{J}}}
\newcommand{\sK}{{\mathscr{K}}}   \newcommand{\sL}{{\mathscr{L}}}
\newcommand{\sM}{{\mathscr{M}}}   \newcommand{\sN}{{\mathscr{N}}}
\newcommand{\sO}{{\mathscr{O}}}   \newcommand{\sP}{{\mathscr{P}}}
\newcommand{\sQ}{{\mathscr{Q}}}   \newcommand{\sR}{{\mathscr{R}}}
\newcommand{\sS}{{\mathscr{S}}}   \newcommand{\sT}{{\mathscr{T}}}
\newcommand{\sU}{{\mathscr{U}}}   \newcommand{\sV}{{\mathscr{V}}}
\newcommand{\sW}{{\mathscr{W}}}   \newcommand{\sX}{{\mathscr{X}}}
\newcommand{\sY}{{\mathscr{Y}}}   \newcommand{\sZ}{{\mathscr{Z}}}


%---------------------------------------
% Math Fraktur Font
%---------------------------------------

%Captital Letters
\newcommand{\mfA}{\mathfrak{A}}	\newcommand{\mfB}{\mathfrak{B}}
\newcommand{\mfC}{\mathfrak{C}}	\newcommand{\mfD}{\mathfrak{D}}
\newcommand{\mfE}{\mathfrak{E}}	\newcommand{\mfF}{\mathfrak{F}}
\newcommand{\mfG}{\mathfrak{G}}	\newcommand{\mfH}{\mathfrak{H}}
\newcommand{\mfI}{\mathfrak{I}}	\newcommand{\mfJ}{\mathfrak{J}}
\newcommand{\mfK}{\mathfrak{K}}	\newcommand{\mfL}{\mathfrak{L}}
\newcommand{\mfM}{\mathfrak{M}}	\newcommand{\mfN}{\mathfrak{N}}
\newcommand{\mfO}{\mathfrak{O}}	\newcommand{\mfP}{\mathfrak{P}}
\newcommand{\mfQ}{\mathfrak{Q}}	\newcommand{\mfR}{\mathfrak{R}}
\newcommand{\mfS}{\mathfrak{S}}	\newcommand{\mfT}{\mathfrak{T}}
\newcommand{\mfU}{\mathfrak{U}}	\newcommand{\mfV}{\mathfrak{V}}
\newcommand{\mfW}{\mathfrak{W}}	\newcommand{\mfX}{\mathfrak{X}}
\newcommand{\mfY}{\mathfrak{Y}}	\newcommand{\mfZ}{\mathfrak{Z}}
%Small Letters
\newcommand{\mfa}{\mathfrak{a}}	\newcommand{\mfb}{\mathfrak{b}}
\newcommand{\mfc}{\mathfrak{c}}	\newcommand{\mfd}{\mathfrak{d}}
\newcommand{\mfe}{\mathfrak{e}}	\newcommand{\mff}{\mathfrak{f}}
\newcommand{\mfg}{\mathfrak{g}}	\newcommand{\mfh}{\mathfrak{h}}
\newcommand{\mfi}{\mathfrak{i}}	\newcommand{\mfj}{\mathfrak{j}}
\newcommand{\mfk}{\mathfrak{k}}	\newcommand{\mfl}{\mathfrak{l}}
\newcommand{\mfm}{\mathfrak{m}}	\newcommand{\mfn}{\mathfrak{n}}
\newcommand{\mfo}{\mathfrak{o}}	\newcommand{\mfp}{\mathfrak{p}}
\newcommand{\mfq}{\mathfrak{q}}	\newcommand{\mfr}{\mathfrak{r}}
\newcommand{\mfs}{\mathfrak{s}}	\newcommand{\mft}{\mathfrak{t}}
\newcommand{\mfu}{\mathfrak{u}}	\newcommand{\mfv}{\mathfrak{v}}
\newcommand{\mfw}{\mathfrak{w}}	\newcommand{\mfx}{\mathfrak{x}}
\newcommand{\mfy}{\mathfrak{y}}	\newcommand{\mfz}{\mathfrak{z}}


\newcommand{\mytitle}{FMA\_3TC13\_TP Poly Exercices 1}
\newcommand{\myauthor}{github/omgeta}
\newcommand{\mydate}{AU 25-26 S1}

\begin{document}
\raggedright
\footnotesize
\begin{center}
{\normalsize{\textbf{\mytitle}}} \\
{\footnotesize{\mydate\hspace{2pt}\textemdash\hspace{2pt}\myauthor}}
\end{center}
\setlist{topsep=-1em, itemsep=-1em, parsep=2em}
\setlist[enumerate]{topsep=0.4em,itemsep=0.8em,parsep=0.35em}
%%%%%%%%%%%%%%%%%%%%%%%%%%%%%%%%%%%%%%%%%%%%%%%%%%%%%%
%                      Begin                         %
%%%%%%%%%%%%%%%%%%%%%%%%%%%%%%%%%%%%%%%%%%%%%%%%%%%%%%
\begin{enumerate}[label=1.\arabic*., start=15]

  \item \begin{minipage}[t]{\linewidth}\raggedright
\textbf{Énoncé.} Soit $f:\RR\to\RR$. Nier les affirmations suivantes de la manière la plus précise possible.

\begin{enumerate}[\arabic*.]
      \item \textbf{Question.} Pour tout $x\in\RR$, $f(x)\leq1$.
\par\smallskip
\textbf{Réponse.} $\exists x \in \RR,\; f(x)>1$.

  \item \textbf{Question.} L'application $f$ est croissante.
\par\smallskip
\textbf{Réponse.} $\exists x,x'\in\RR,\;x<x' \text{ et } f(x)>f(x')$.

  \item \textbf{Question.} L'application $f$ est croissante et positive.
\par\smallskip
\textbf{Réponse.} $(\exists x<x',\;f(x)>f(x'))\lor(\exists x\in\RR,\;f(x)<0)$.

  \item \textbf{Question.} Il existe $x\in\RR_{\geq0}$ tel que $f(x)\leq0$.
\par\smallskip
\textbf{Réponse.} $\forall x\in\RR_{\geq0},\;f(x)>0$.

\end{enumerate}
\end{minipage}

  \item \begin{minipage}[t]{\linewidth}\raggedright
\textbf{Énoncé.} Pour chaque fonction, dire si elle est injective, surjective, bijective. Déterminer $f^{-1}(\{0\})$, $f^{-1}(\{0,1\})$, $f^{-1}(\NN)$ et l'image de $f$.

\begin{enumerate}[\arabic*.]
      \item \textbf{Question.} $f_1:\ZZ\to\ZZ$, $n\mapsto2n$.
\par\smallskip
\textbf{Réponse.} $f_1$ est injective, mais non surjective : $2n=2m\Rightarrow n=m$, et $1$ n'a pas d'antécédent.\\
      $f_1^{-1}(\{0\})=\{0\}$, $f_1^{-1}(\{0,1\})=\{0\}$, $f_1^{-1}(\NN)=\NN$, $f_1(\ZZ)=2\ZZ$.

  \item \textbf{Question.} $f_2:\ZZ\to\ZZ$, $n\mapsto-n$.
\par\smallskip
\textbf{Réponse.} $f_2$ est bijective car $f_2\circ f_2=\operatorname{id}_{\ZZ}$.\\
      $f_2^{-1}(\{0\})=\{0\}$, $f_2^{-1}(\{0,1\})=\{0,-1\}$, $f_2^{-1}(\NN)=-\NN$, $f_2(\ZZ)=\ZZ$.

  \item \textbf{Question.} $f_3:\RR\to\RR$, $x\mapsto x^2$.
\par\smallskip
\textbf{Réponse.} $f_3$ n'est ni injective ($f_3(1)=f_3(-1)$), ni surjective ($-1$ n'a pas d'antécédent).\\
      $f_3^{-1}(\{0\})=\{0\}$, $f_3^{-1}(\{0,1\})=\{-1,0,1\}$,\\
      $f_3^{-1}(\NN)=\{\pm\sqrt n:n\in\NN\}$, $f_3(\RR)=\RR_{\geq0}$.

  \item \textbf{Question.} $f_4:\RR\to\RR_{\geq0}$, $x\mapsto x^2$.
\par\smallskip
\textbf{Réponse.} $f_4$ est surjective car $y=f_4(\sqrt y)$ pour tout $y\geq0$, mais non injective.\\
      $f_4^{-1}(\{0\})=\{0\}$, $f_4^{-1}(\{0,1\})=\{-1,0,1\}$,\\
      $f_4^{-1}(\NN)=\{\pm\sqrt n:n\in\NN\}$, $f_4(\RR)=\RR_{\geq0}$.

  \item \textbf{Question.} $f_5:\CC\to\CC$, $z\mapsto z^2$.
\par\smallskip
\textbf{Réponse.} $f_5$ est surjective car tout complexe possède une racine carrée, mais non injective.\\
      $f_5^{-1}(\{0\})=\{0\}$, $f_5^{-1}(\{0,1\})=\{-1,0,1\}$,\\
      $f_5^{-1}(\NN)=\{\pm\sqrt n:n\in\NN\}$, $f_5(\CC)=\CC$.

\end{enumerate}
\end{minipage}

  \item \begin{minipage}[t]{\linewidth}\raggedright
\textbf{Énoncé.} Soient $f:A\to B$, $g:B\to C$, $h:C\to D$. Montrer les implications suivantes.

\begin{enumerate}[\arabic*.]
      \item \textbf{Question.} Si $g\circ f$ est injective, alors $f$ est injective.
\par\smallskip
\textbf{Réponse.} Si $g\circ f$ est injective, alors $f$ est injective
      \begin{enumproof}
        \item Soient $x,x'\in A$ tels que $f(x)=f(x')$
        \item Alors $g(f(x))=g(f(x'))$
        \item Donc $x=x'$\hfill(Injectivité de $g\circ f$)
        \item $\therefore f$ est injective
      \end{enumproof}

  \item \textbf{Question.} Si $g\circ f$ est surjective, alors $g$ est surjective.
\par\smallskip
\textbf{Réponse.} Si $g\circ f$ est surjective, alors $g$ est surjective
      \begin{enumproof}
        \item Soit $z\in C$. Il existe $x\in A$ tel que $g(f(x))=z$
        \item Alors $y=f(x)\in B$ vérifie $g(y)=z$
        \item $\therefore g$ est surjective
      \end{enumproof}

  \item \textbf{Question.} Si $g\circ f$ et $h\circ g$ sont bijectives, alors $f,g,h$ sont bijectives.
\par\smallskip
\textbf{Réponse.} Si $g\circ f$ et $h\circ g$ sont bijectives, alors $f,g,h$ sont bijectives
      \begin{enumproof}
        \item $h\circ g$ injective implique $g$ injective\hfill(Question 1)
        \item $g\circ f$ surjective implique $g$ surjective\hfill(Question 2)
        \item Donc $g$ est bijective
        \item $f=g^{-1}\circ(g\circ f)$ et $h=(h\circ g)\circ g^{-1}$ sont des compositions de bijections
        \item $\therefore f,g,h$ sont bijectives
      \end{enumproof}

\end{enumerate}
\end{minipage}

  \item \begin{minipage}[t]{\linewidth}\raggedright
\textbf{Énoncé.} Soit $f:X\to Y$. Sur les ensembles de parties, on définit $\hat f:\mathcal P(X)\to\mathcal P(Y)$, $A\mapsto f(A)$, et $\check f:\mathcal P(Y)\to\mathcal P(X)$, $B\mapsto f^{-1}(B)$.

\begin{enumerate}[\arabic*.]
      \item \textbf{Question.} Montrer que $f$ est injective si et seulement si $\hat f$ est injective.
\par\smallskip
\textbf{Réponse.} $f$ est injective si et seulement si $\hat f$ est injective
      \begin{enumproof}
        \item ($\Rightarrow$) Supposons $f$ injective et $f(A)=f(A')$.
          Pour $x\in A$, il existe $x'\in A'$ tel que $f(x)=f(x')$.
          Alors $x=x'\in A'$, donc $A\subseteq A'$.
          De même $A'\subseteq A$, donc $A=A'$.
        \item ($\Leftarrow$) Si $f(x)=f(x')$, alors $\hat f(\{x\})=\hat f(\{x'\})$.
          L'injectivité de $\hat f$ donne $\{x\}=\{x'\}$, donc $x=x'$.
      \end{enumproof}

  \item \textbf{Question.} Montrer que $f$ est surjective si et seulement si $\check f$ est injective.
\par\smallskip
\textbf{Réponse.} $f$ est surjective si et seulement si $\check f$ est injective
      \begin{enumproof}
        \item ($\Rightarrow$) Supposons $f$ surjective et $f^{-1}(B)=f^{-1}(B')$.
          Si $y\in B$, choisissons $x\in X$ tel que $f(x)=y$.
          Alors $x\in f^{-1}(B)=f^{-1}(B')$, donc $y\in B'$.
          Par symétrie, $B=B'$.
        \item ($\Leftarrow$) Si $f$ n'est pas surjective, il existe $y\in Y$ sans antécédent.
          Alors $\check f(\{y\})=\varnothing=\check f(\varnothing)$, alors que $\{y\}\neq\varnothing$.
          Ceci contredit l'injectivité de $\check f$.
      \end{enumproof}

\end{enumerate}
\end{minipage}

\end{enumerate}
%%%%%%%%%%%%%%%%%%%%%%%%%%%%%%%%%%%%%%%%%%%%%%%%%%%%%%
%                       End                          %
%%%%%%%%%%%%%%%%%%%%%%%%%%%%%%%%%%%%%%%%%%%%%%%%%%%%%%

\end{document}
